% !PRINTED 528
\weberpage{528}
From $P_2$ likewise three conjugate groups arise, which however have no common divisor other than 1.

\gesperrt{Thus, apart from $P$, $Q$, $P_1$, $Q_1$, $P_2$ and the divisors of $P$ conjugate with $P_1$ and $P_2$, there are no transitive permutation groups of four elements.}

Among the intransitive groups, the group (3):
\[
P_3 = 1, \ (0,2), \ (1,3), \ (0,2)(1,3),
\]
still deserves to be singled out, which likewise gives rise to a system of three conjugate groups.

\begin{center}
\S.\,161.

\gesperrt{Solution of the biquadratic equations.}
\end{center}

The various methods of solution of the biquadratic equations differ from one another in the order in which the various divisors of the group $P$ are employed.

Let $\alpha$, $\alpha_1$, $\alpha_2$, $\alpha_3$ be the roots of the biquadratic equation
\weq{1}{f(x) = x^4 - a_1x^3 + a_2x^2 - a_3x + a_4 = 0.}
We denote the discriminant by $D$ and set
\weq{2}{\sqrt{D} = (\alpha-\alpha_1)(\alpha-\alpha_2)(\alpha-\alpha_3)(\alpha_1-\alpha_2)(\alpha_1-\alpha_3)(\alpha_2-\alpha_3).}
Moreover we recall the expressions for the invariants (\S.\,64):
\weq{3}{A = a_2^2 - 3a_1a_3 + 12a_4}
\weq{4}{B = 27a_1^2a_4 + 27a_3^2 + 2a_2^3 - 72a_2a_4 - 9a_1a_2a_3,}
by means of which $D$ is expressed in the form:
\weq{5}{27D = 4A^3 - B^2.}
If one first adjoins $\sqrt{D}$, the group of the equation reduces to $Q$, and if one then seeks a function belonging to $Q_1$, the group is reduced by it to $Q_1$. Since $Q_1$ has index 3 within $Q$, such a function will depend on a cubic equation, and indeed, since $Q_1$ is a normal divisor of $Q$, on a normal equation. Let us take, for example,
\weq{6}{y = (\alpha-\alpha_1)(\alpha_2-\alpha_3)}
%CONT>>
